% Chapter 1 -- The Real Field, 1.19-1.22
\rsection{The Real Field}

\begin{signpost}[The hinge of the chapter]
Everything so far has been preparation, and everything after is consequence.
This section contains the one theorem the chapter was written to state.

\smallskip
Watch what Rudin does with it. He states 1.19, declines to prove it, and then
immediately proves 1.20 and 1.21 --- results he could have extracted from the
construction ``with very little extra effort,'' as he says. He derives them from
the least-upper-bound property instead, on purpose: he wants you to see the
property doing work before you see where it came from. Read 1.20 and 1.21 less
as facts about $\R$ than as two worked demonstrations of how a supremum is
used.
\end{signpost}

\noindent
We now state the existence theorem which is the core of this chapter.

\ritem{1.19}{Theorem}
\emph{There exists an ordered field $\R$ which has the least-upper-bound
property.}

\emph{Moreover, $\R$ contains $\Q$ as a subfield.}

The second statement means that $\Q \subset \R$ and that the operations of
addition and multiplication in $\R$, when applied to members of $\Q$, coincide
with the usual operations on rational numbers; also, the positive rational
numbers are positive elements of $\R$.\kn{``Subfield'' is a stronger claim than
``subset.'' It says the copy of $\Q$ sitting inside $\R$ has the arithmetic and
the order it had before --- so nothing you already know about rationals is
invalidated by the enlargement.} The members of $\R$ are called \emph{real
numbers}.

The proof of Theorem 1.19 is rather long and a bit tedious and is therefore
presented in an Appendix to Chap.\ 1. The proof actually constructs $\R$ from
$\Q$.\kd{Worth pausing on. Everything in the next ten chapters is a consequence
of this one theorem plus the field and order axioms. If you accept 1.19, you
never need the appendix; if you doubt it, the appendix settles it by exhibiting
$\R$ as the set of Dedekind cuts of $\Q$. Either way the rest of the book is
unaffected, which is why Rudin can safely postpone it.}

The next theorem could be extracted from this construction with very little
extra effort. However, we prefer to derive it from Theorem 1.19 since this
provides a good illustration of what one can do with the least-upper-bound
property.

\ritem{1.20}{Theorem}
\begin{enumerate}[label=(\alph*), leftmargin=2.2em, itemsep=3pt, topsep=3pt]
  \item \emph{If $x \in \R$, $y \in \R$, and $x > 0$, then there is a positive
        integer $n$ such that}
        \[ nx > y. \]
  \item \emph{If $x \in \R$, $y \in \R$, and $x < y$, then there exists a
        $p \in \Q$ such that $x < p < y$.}
\end{enumerate}
Part (a) is usually referred to as the \emph{archimedean property} of $\R$. Part
(b) may be stated by saying that $\Q$ is \emph{dense} in $\R$: Between any two
real numbers there is a rational one.\kw{Part (a) is a theorem, not an axiom,
and it is not automatic. There are ordered fields in which it fails --- fields
containing an element larger than every integer. Such fields cannot have the
least-upper-bound property, and (a) is precisely the proof of that.}

\begin{proof}
(a) Let $A$ be the set of all $nx$, where $n$ runs through the positive
integers. If (a) were false, then $y$ would be an upper bound of $A$. But then
$A$ has a least upper bound in $\R$. Put $\alpha = \sup A$. Since $x > 0$,
$\alpha - x < \alpha$, and $\alpha - x$ is not an upper bound of $A$. Hence
$\alpha - x < mx$ for some positive integer $m$. But then
$\alpha < (m+1)x \in A$, which is impossible, since $\alpha$ is an upper bound
of $A$.

(b) Since $x < y$, we have $y - x > 0$, and (a) furnishes a positive integer $n$
such that
\[ n(y-x) > 1. \]
Apply (a) again, to obtain positive integers $m_1$ and $m_2$ such that
$m_1 > nx$, $m_2 > -nx$. Then
\[ -m_2 < nx < m_1. \]
Hence there is an integer $m$ (with $-m_2 \le m \le m_1$) such that
\[ m - 1 \le nx < m. \]
If we combine these inequalities, we obtain
\[ nx < m \le 1 + nx < ny. \]
Since $n > 0$, it follows that
\[ x < \frac{m}{n} < y. \]
This proves (b), with $p = m/n$.
\end{proof}

\begin{unpack}[Part (a): grab the only handle]
How would one find this proof? Negate the conclusion: $nx \le y$ for every
$n$. That is a \emph{boundedness} statement, and the only machinery 1.19
offers is that bounded sets have suprema --- so the proof is forced to form
$A = \{nx\}$ and take $\alpha = \sup A$. There is nothing else to try.

The attack on $\alpha$ is then the standard one from 1.8(ii): step back by a
positive amount --- here $x$ itself, the natural unit --- to find some
$mx > \alpha - x$; step forward one term to $(m+1)x > \alpha$, and the
supremum is beaten by a member of its own set. Keep the pattern: \emph{when
the hypothesis is ``bounded,'' the proof is ``take the sup and derive a
contradiction one step past it.''} It reappears at 1.21, 2.28, 2.38 and 3.14,
and it is pattern 3 of the chapter's strategy index.
\end{unpack}

\begin{unpack}[The three moves in part (b)]
Rudin gives the computation and not the plan. The plan is this: to fit a
rational strictly between $x$ and $y$, first \emph{stretch} the gap until it is
wider than $1$, then \emph{catch} an integer inside the stretched gap, then
\emph{shrink} back.

\smallskip
\textbf{Stretch.} Part (a) with the roles $x \leadsto y-x$ and $y \leadsto 1$
gives an $n$ with $n(y-x) > 1$. The interval $(nx, ny)$ now has length greater
than $1$.

\smallskip
\textbf{Catch.} An interval of length greater than $1$ must contain an integer.
Rudin proves this rather than asserting it, and that is what $m_1, m_2$ are for:
they trap $nx$ between two integers, so the set of integers exceeding $nx$ is
nonempty and bounded below, and therefore has a least element $m$. That gives
$m-1 \le nx < m$. The existence of a \emph{least} such integer is well-ordering
again --- the same principle used silently in 1.1. Rudin never states it as an
axiom, but this step and the archimedean property both consume it.

\smallskip
\textbf{Shrink.} From $m \le 1 + nx$ and $1 < n(y-x)$ we get
$m < nx + n(y-x) = ny$, and dividing by $n > 0$ puts $m/n$ strictly between $x$
and $y$.

\medskip
\begin{center}
\begin{tikzpicture}[x=1mm, y=1mm]
  % ---- stage 1: the original gap, possibly tiny
  \draw[axis] (-1,26) -- (56,26);
  \node[ptfill, minimum size=2.6pt] at (18,26) {};
  \node[ptfill, minimum size=2.6pt] at (25,26) {};
  \node[mlbl, anchor=north] at (18,24.5) {$x$};
  \node[mlbl, anchor=north] at (25,24.5) {$y$};
  \node[lbl, anchor=west] at (60,26) {the gap may be tiny};
  \node[lbl, anchor=east] at (-3,26) {\textbf{1}};

  % ---- stage 2: stretched by n, now longer than 1
  \draw[axis] (-1,13) -- (56,13);
  \foreach \i in {6,12,18,24,30,36,42,48}{\draw[thin] (\i,11.7) -- (\i,14.3);}
  \node[ptfill, minimum size=2.6pt] at (15,13) {};
  \node[ptfill, minimum size=2.6pt] at (45,13) {};
  \node[mlbl, anchor=north] at (15,11.3) {$nx$};
  \node[mlbl, anchor=north] at (45,11.3) {$ny$};
  \node[ptopen] at (18,13) {};
  \node[mlbl, anchor=south] at (18,14.8) {$m$};
  \node[lbl, anchor=west] at (60,15) {length $>1$, so an};
  \node[lbl, anchor=west] at (60,11.5) {integer is trapped};
  \node[lbl, anchor=east] at (-3,13) {\textbf{2}};

  % ---- stage 3: divided back by n
  \draw[axis] (-1,0) -- (56,0);
  \node[ptfill, minimum size=2.6pt] at (18,0) {};
  \node[ptfill, minimum size=2.6pt] at (25,0) {};
  \node[ptopen] at (20,0) {};
  \node[mlbl, anchor=north] at (18,-1.6) {$x$};
  \node[mlbl, anchor=north] at (25,-1.6) {$y$};
  \node[mlbl, anchor=south] at (20,1.8) {$m/n$};
  \node[lbl, anchor=west] at (60,0) {divide back by $n$};
  \node[lbl, anchor=east] at (-3,0) {\textbf{3}};

  \draw[guide, -{Stealth[length=3.5pt]}] (33,23) -- (33,17);
  \draw[guide, -{Stealth[length=3.5pt]}] (33,10) -- (33,4);
\end{tikzpicture}
\end{center}
\figcap{Stretch until an integer must fall inside, catch it, then shrink back.
The integer $m$ becomes the rational $m/n$ between $x$ and $y$.}
\end{unpack}

\medskip
\noindent
We shall now prove the existence of $n$th roots of positive reals. This proof
will show how the difficulty pointed out in the Introduction (irrationality of
$\sqrt 2$) can be handled in $\R$.

\ritem{1.21}{Theorem}
\emph{For every real $x>0$ and every integer $n>0$ there is one and only one
positive real $y$ such that $y^n = x$.}

This number $y$ is written $\sqrt[n]{x}$ or $x^{1/n}$.

\begin{proof}
That there is at most one such $y$ is clear, since $0 < y_1 < y_2$ implies
$y_1^n < y_2^n$.

Let $E$ be the set consisting of all positive real numbers $t$ such that
$t^n < x$. If $t = x/(1+x)$ then $0 \le t < 1$. Hence $t^n \le t < x$. Thus
$t \in E$, and $E$ is not empty.\kn{Two hypotheses of the \lub{} property are
being discharged, and Rudin spends a sentence on each. This paragraph is
nonemptiness; the next is boundedness. Both are required before $\sup E$ may be
written down --- compare the warning at 1.10.}

If $t > 1+x$ then $t^n \ge t > x$, so that $t \notin E$. Thus $1+x$ is an upper
bound of $E$.

Hence Theorem 1.19 implies the existence of
\[ y = \sup E. \]

To prove that $y^n = x$ we will show that each of the inequalities $y^n < x$ and
$y^n > x$ leads to a contradiction.\kd{Note which half of Definition 1.8 each
case attacks. Assuming $y^n < x$ will contradict ``$y$ is an \emph{upper
bound}''; assuming $y^n > x$ will contradict ``$y$ is the \emph{least} upper
bound.'' Both clauses get used exactly once, which is a reliable sign that $E$
was chosen well.}

The identity $b^n - a^n = (b-a)(b^{n-1} + b^{n-2}a + \cdots + a^{n-1})$ yields
the inequality
\[ b^n - a^n < (b-a)nb^{n-1} \]
when $0 < a < b$.

\begin{center}
\begin{tikzpicture}[x=1mm, y=1mm]
  \draw[axis] (-3,0) -- (95,0);
  \foreach \x in {0,19,47.5,90.25}{\draw[thin] (\x,-1.6) -- (\x,1.6);}
  \node[mlbl, anchor=north] at (0,-2.6)    {$a^3$};
  \node[mlbl, anchor=north] at (19,-2.6)   {$a^2b$};
  \node[mlbl, anchor=north] at (47.5,-2.6) {$ab^2$};
  \node[mlbl, anchor=north] at (90.25,-2.6){$b^3$};
  \draw[thin, {Stealth[length=2.5pt]}-{Stealth[length=2.5pt]}]
    (0.6,4) -- (18.4,4);
  \draw[thin, {Stealth[length=2.5pt]}-{Stealth[length=2.5pt]}]
    (19.6,4) -- (46.9,4);
  \draw[thin, {Stealth[length=2.5pt]}-{Stealth[length=2.5pt]}]
    (48.1,4) -- (89.65,4);
  \node[mlbl, anchor=south] at (9.5,5)   {$a^2(b{-}a)$};
  \node[mlbl, anchor=south] at (33.25,5) {$ab(b{-}a)$};
  \node[mlbl, anchor=south] at (68.9,5)  {$b^2(b{-}a)$};
\end{tikzpicture}
\end{center}
\figcap{The identity, drawn for $n = 3$: the walk from $a\!\cdot\!a\!\cdot\!a$
to $b\!\cdot\!b\!\cdot\!b$ replaces one factor at a time, cutting $[a^3, b^3]$
into three pieces whose widths are the three terms of the identity. Since
$a<b$, every piece is shorter than the last one, whence
$b^n - a^n < (b-a)nb^{n-1}$. (Drawn for $a=1$, $b=1.5$. For $n=2$ this is the
difference of squares; at 3.26 the same identity, with $b=1$, $a=x$, sums the
geometric series.)}

Assume $y^n < x$. Choose $h$ so that $0 < h < 1$ and
\[ h < \frac{x-y^n}{n(y+1)^{n-1}}. \]
Put $a = y$, $b = y+h$. Then
\[ (y+h)^n - y^n < hn(y+h)^{n-1} < hn(y+1)^{n-1} < x - y^n. \]
Thus $(y+h)^n < x$ and $y+h \in E$. Since $y+h > y$, this contradicts the fact
that $y$ is an upper bound of $E$.\kn{Do not try to see the choice of $h$ at a
glance --- nobody does. It was found \emph{backwards}: write the goal
$(y+h)^n < x$, bound the left side by the displayed inequality, and solve for
$h$; the strange side condition $h<1$ exists only to get $h$ out of the
exponent first. The box after the proof reconstructs the whole derivation, and
the $k$ of the next paragraph is found the same way.}

Assume $y^n > x$. Put
\[ k = \frac{y^n - x}{ny^{n-1}}. \]
Then $0 < k < y$. If $t \ge y-k$, we conclude that
\[ y^n - t^n \le y^n - (y-k)^n < kny^{n-1} = y^n - x. \]
Thus $t^n > x$ and $t \notin E$. It follows that $y - k$ is an upper bound of
$E$. But $y - k < y$, which contradicts the fact that $y$ is the least upper
bound of $E$.

Hence $y^n = x$, and the proof is complete.
\end{proof}

\begin{proofwalk}[How the proof flows]
  \item \textbf{Choose the method.} A number with $y^n = x$ must be
        \emph{produced}, and the only production tool is the least-upper-bound
        property --- pattern 3 again. So the proof must design a set whose
        supremum is the root. The natural candidate: everything that
        undershoots, $E = \{t>0 : t^n < x\}$.
  \item \textbf{Buy the supremum.} Both tickets, each bought with one
        sentence: $x/(1+x)$ lies in $E$ (nonempty), $1+x$ bounds it. Rudin's
        care here is the warning of 1.10 being heeded, and $y = \sup E$
        exists.
  \item \textbf{Set up the pincer.} By trichotomy (1.5), it suffices to
        kill $y^n < x$ and $y^n > x$. There is no direct computation of
        $y^n$ available --- $y$ is only known as a supremum --- so each case
        must be turned into a statement \emph{about bounds}, the only
        language $y$ speaks.
  \item \textbf{Kill $y^n < x$.} If $y$ undershoots, continuity-style
        estimates (the binomial identity) produce a slightly larger
        $y+h$ still in $E$ --- contradicting clause (i) of 1.8: $y$ was
        supposed to be an upper bound. The choice of $h$ is reverse-engineered
        in the next box.
  \item \textbf{Kill $y^n > x$.} If $y$ overshoots, the same estimate run
        downward shows $y-k$ already bounds $E$ --- contradicting clause (ii):
        $y$ was supposed to be the \emph{least} upper bound.
  \item \textbf{Audit.} Clause (i) died in step 4, clause (ii) in step 5 ---
        each spent exactly once, as in 1.11. When both halves of the
        supremum's definition are consumed, the set $E$ was the right one;
        when a proof of this shape never uses one clause, suspect the set.
\end{proofwalk}

\begin{center}
\begin{tikzpicture}[x=1mm, y=1mm]
  \node[rmbox, text width=30mm] (e) at (0,40)
    {$E = \{t>0 : t^n<x\}$\\nonempty ($t = \frac{x}{1+x}$),\\bounded (by $1+x$)};
  \node[rmbox, text width=24mm] (y) at (44,40)
    {$y = \sup E$\\exists};
  \draw[rmarrow] (e) -- node[lbl, anchor=south, inner sep=2pt]
    {1.19} (y);
  \node[rmbox, text width=35mm] (c1) at (14,18)
    {if $y^n < x$: nudge up,\\$y+h \in E$ ---\\$y$ was no upper bound};
  \node[rmbox, text width=35mm] (c2) at (60,18)
    {if $y^n > x$: step down,\\$y-k$ still bounds $E$ ---\\$y$ was not the least};
  \draw[rmarrow] (y) -- node[lbl, anchor=east, inner sep=2pt, pos=0.72]
    {contra 1.8(i)} (c1);
  \draw[rmarrow] (y) -- node[lbl, anchor=west, inner sep=2pt, pos=0.72]
    {contra 1.8(ii)} (c2);
  \node[rmbox, text width=22mm] (c) at (38,-3) {$y^n = x$};
  \draw[rmarrow] (c1) -- (c);
  \draw[rmarrow] (c2) -- (c);
\end{tikzpicture}
\end{center}
\figcap{The proof as a flow diagram: qualify $E$, mint $y = \sup E$, then kill
both strict inequalities --- each contradiction consuming a different clause of
Definition 1.8. Trichotomy (1.5) does the rest.}

\begin{unpack}[Reverse-engineering the choice of $h$]
Nobody guesses that $h$. It is found by writing the goal down and solving
backwards, and Rudin has simply deleted the derivation.

\smallskip
\textbf{Goal.} Find $h>0$ with $(y+h)^n < x$; equivalently
$(y+h)^n - y^n < x - y^n$, whose right-hand side is a fixed positive number
because we assumed $y^n < x$.

\smallskip
\textbf{Bound the left side.} We cannot compute $(y+h)^n - y^n$, but the
displayed inequality with $a = y$, $b = y+h$ bounds it by $hn(y+h)^{n-1}$. So it
suffices to make $hn(y+h)^{n-1} \le x - y^n$.

\smallskip
\textbf{Get $h$ out of the exponent.} That inequality still has $h$ buried
inside $(y+h)^{n-1}$ and cannot be solved as it stands. Kill it by agreeing in
advance that $h < 1$; then $(y+h)^{n-1} < (y+1)^{n-1}$, a constant. This is the
\emph{only} reason the condition $0 < h < 1$ appears --- it is bookkeeping, not
mathematics.

\smallskip
\textbf{Solve.} It now suffices that $hn(y+1)^{n-1} < x - y^n$, that is,
$h < (x-y^n)/\bigl(n(y+1)^{n-1}\bigr)$, which is Rudin's line. Both conditions
on $h$ are upper bounds, so any $h$ below the smaller of the two works, and such
an $h$ exists.

\smallskip
The second case is the same argument run downward, and $k$ is found the same
way; there the constraint $k < y$ plays the role that $h < 1$ played here.

\medskip
\begin{center}
\begin{tikzpicture}[x=1mm, y=1mm]
  \draw[axis] (-1,0) -- (94,0);
  \draw[seg] (6,0) -- (48,0);
  \node[ptopen] at (48,0) {};
  \node[lbl, anchor=north] at (24,-2) {$E=\{t>0 : t^n<x\}$};
  \draw[guide] (48,-7) -- (48,20);
  \node[lbl, anchor=south] at (48,20) {$y=\sup E$};

  % case y^n < x : nudge up and stay inside E
  \node[ptfill, minimum size=2.6pt] at (58,8) {};
  \draw[thin, -{Stealth[length=3.5pt]}] (48,8) -- (58,8);
  \node[mlbl, anchor=south] at (53,9) {$h$};
  \node[lbl, anchor=west] at (60,8) {if $y^n<x$: $y+h$ still in $E$,};
  \node[lbl, anchor=west] at (60,4.6) {so $y$ was not an upper bound};

  % case y^n > x : step down and still bound
  \node[ptfill, minimum size=2.6pt] at (38,-11) {};
  \draw[thin, -{Stealth[length=3.5pt]}] (48,-11) -- (38,-11);
  \node[mlbl, anchor=south] at (43,-10) {$k$};
  \node[lbl, anchor=west] at (52,-11) {if $y^n>x$: $y-k$ still bounds $E$,};
  \node[lbl, anchor=west] at (52,-14.4) {so $y$ was not the \emph{least} one};
\end{tikzpicture}
\end{center}
\figcap{Both halves of Definition 1.8 get used exactly once --- a reliable sign
that $E$ was the right set to take a supremum of.}
\end{unpack}

\begin{deeper}[What the proof is really doing]
Strip the algebra and the argument reads: \emph{$t \mapsto t^n$ is continuous,
so if $y^n$ were strictly below $x$ we could nudge $y$ upward and stay below
$x$.} The choice of $h$ is a hand-rolled continuity estimate, written out
because continuity is not available until 4.4 --- and the machinery of Chapter 4
is not available until $\R$ exists, which is what we are in the middle of
establishing.

Once Chapter 4 is in place the same theorem falls out of the intermediate value
theorem (4.23) in three lines. Rudin proves it here anyway because he needs
$n$th roots long before then, and because it is the first genuine demonstration
that the \lub{} property produces numbers $\Q$ could not.
\end{deeper}

\begin{pitfall}[Why $E$ is restricted to positive $t$]
Without positivity the uniqueness claim is simply false: for $n = 2$, $x = 4$,
both $2$ and $-2$ are roots. Rudin's ``one and only one \emph{positive} real''
is doing real work, and the restriction on $E$ is what licenses the opening line
of the proof, $0 < y_1 < y_2 \Rightarrow y_1^n < y_2^n$, which requires both
numbers to be positive.
\end{pitfall}

\ritem{}{Corollary}
\emph{If $a$ and $b$ are positive real numbers and $n$ is a positive integer,
then}
\[ (ab)^{1/n} = a^{1/n}b^{1/n}. \]

\begin{proof}
Put $\alpha = a^{1/n}$, $\beta = b^{1/n}$. Then
\[ ab = \alpha^n\beta^n = (\alpha\beta)^n \]
since multiplication is commutative. [Axiom (M2) in Definition 1.12.] The
uniqueness assertion of Theorem 1.21 shows therefore that
\[ (ab)^{1/n} = \alpha\beta = a^{1/n}b^{1/n}. \qedhere \]
\end{proof}

\begin{unpack}[Uniqueness doing the work]
Note what carries the argument: not a computation with roots, but the
\emph{uniqueness} half of 1.21. We exhibit one positive number, namely
$\alpha\beta$, whose $n$th power is $ab$; since 1.21 says there is only one such
number, that number must be $(ab)^{1/n}$.

This is the standard way to prove identities about $n$th roots, and it is worth
internalising because the direct approach is hopeless --- there is no formula
for $a^{1/n}$ to manipulate. Rudin uses the same move at 1.33(c) to get
$|zw| = |z||w|$.
\end{unpack}

\ritem{1.22}{Decimals}
We conclude this section by pointing out the relation between real numbers and
decimals.

Let $x > 0$ be real. Let $n_0$ be the largest integer such that $n_0 \le x$.
(Note that the existence of $n_0$ depends on the archimedean property of
$\R$.)\kn{Rudin flags the dependence himself, which is unusual for him and a
signal that it matters. Without 1.20(a) there might be no integer at all
below --- or above --- a given $x$, and the decimal expansion could not even
start.} Having chosen $n_0, n_1, \dots, n_{k-1}$, let $n_k$ denote the largest
integer such that
\[ n_0 + \frac{n_1}{10} + \cdots + \frac{n_k}{10^k} \le x. \]
Let $E$ be the set of these numbers
\begin{equation}
  n_0 + \frac{n_1}{10} + \cdots + \frac{n_k}{10^k}
  \qquad (k = 0,1,2,\dots). \tag{1.5}
\end{equation}
Then $x = \sup E$. The decimal expansion of $x$ is
\begin{equation}
  n_0.n_1n_2n_3\cdots. \tag{1.6}
\end{equation}
Conversely, for any infinite decimal expansion (1.6) the set $E$ of numbers
(1.5) is bounded above and (1.6) is the decimal expansion of $\sup E$.

Since we shall never use decimals, we do not enter into a detailed
discussion.\kw{Take that last sentence seriously --- it is a warning, not
modesty. Decimals are a \emph{representation} of real numbers, not a definition
of them, and reasoning about $\R$ through decimal expansions is a reliable way
to go wrong (the two expansions of $1 = 0.999\dots$ are the familiar symptom).
Everything in this book is proved from 1.19 and never from digits.}

\begin{deeper}[Why this passage is here at all]
It answers the question the Introduction raised and then dropped: the sequence
$1, 1.4, 1.41, \dots$ ``tends to $\sqrt2$'' --- what does it tend to? The answer
is now available and is stated in one word: $\sup$. The truncations form a set
$E$, that set is bounded above, and 1.19 hands you its supremum. The number the
decimals were approximating is that supremum, and nothing more mysterious.

Note the two directions Rudin states. Every real has a decimal expansion, and
every decimal expansion names a real. Together they say the informal picture of
$\R$ you arrived with is not wrong --- it just could not be a foundation, since
$E$ having a supremum is exactly what $\Q$ failed to provide.
\end{deeper}
